\documentclass[12pt]{article}

\usepackage{helpers}

\begin{document}

Name: \makebox[3in]{\hrulefill
%NAME
\hrulefill}

\vfill

\begin{center}
{\huge Test 2}

{\large Due March 27, 2026}
\end{center}

\testrules

\vfill

\pledge

\newpage

\section*{Assembly Programming}

Demonstrate understanding of assembly programming by writing a program using global variables, local variables, input with \texttt{scanf}, and/or output with \texttt{printf}.\\

%BEGIN_STANDARD_AP

Write an assembly program called \texttt{square.s} that does the following:
\begin{itemize}
\item Ask the user to enter a number, and store their response in a \textbf{global variable} called \texttt{num}.
\item Compute the square of \texttt{num}, by loading it's value from memory, multiplying it by itself, and storing the result to a \textbf{local variable} called \texttt{square}.
\item Print the sentence ``The square of \{num\} is \{square\}'', where \{num\} is replaced with the value of the global variable \texttt{num}, and \{square\} is replaced with the value of the local variable \texttt{square}. For example, if the user entered 7, the program would print ``The square of 7 is 49''.
\end{itemize}

The next two pages outline the structure of the program with comments. You may write your assembly code between the comments, or on a separate sheet of paper. If you write on a separate sheet of paper, you do not need to rewrite the comments.\\

\newpage

\textbf{square.s}

\begin{verbatim}
@ Global constants
@ prompt: string asking the user to enter a number



@ input_string: formatting string for reading the user's input



@ print_report: formatting string for reporting the square of the user's number



@ Global variables
@ num: integer, the number entered by the user



@ main
     @ set up stack frame for main, one local variable



     @ ask the user to enter a number


     
     @ store the response in num


     
     @ load the value of num from memory



     @ multiple the value by itself, store the result in square


     
     @ load the values of num and square from memory


     
     @ print a sentence saying what the square of number is


     
     @ return 0, stack frame teardown


\end{verbatim}



\vfill

\standardsfooter

%END_STANDARD_AP

\newpage

\section*{Hardware Components}

Demonstrate understanding of hardware components and their purposes through answering conceptual questions and identifying components.\\

%BEGIN_STANDARD_HC

\begin{enumerate}
\item What is the role of the CPU in a computer?
\vfill
\item What is the role of main memory in a computer?
\vfill
\item What is the role of the control unit in a computer?
\vfill
\item What is the role of the processing unit in a computer? 
\vfill
\item If the clock rate is 4 GHz, how long does it take one instruction to complete? Explain your answer.
\vfill
\item If the clock rate is 1 GHz, how many ticks per second does the computer have? Explain your answer.
\vfill
\end{enumerate}

\pagebreak

For each of the following, include which hardware components are involved, information sent on the address, control, and/or data bus, and the use of and updates to special registers.

\begin{enumerate}

  \setcounter{enumi}{6}

\item Write a detailed description of what happens in the fetch phase of instruction execution in the von Neumann architecture.
\vfill
\item Write a detailed description of what happens in the decode phase of instruction execution in the von Neumann architecture.
\vfill
\item Write a detailed description of what happens in the execute phase of instruction execution in the von Neumann architecture.
\vfill
\item Write a detailed description of what happens in the store phase of instruction execution in the von Neumann architecture.
\vfill
\end{enumerate}


\vfill

\standardsfooter

%END_STANDARD_HC

\end{document}